\subsection{双线性映射。}

本小节中，A 与 B 表示两个环，E 表示左 A-模，F 表示右 B-模，G 表示 (A, B)-双模，即一个交换群，配备左 A-模结构与右 B-模结构，使得 \((ag)b = a(gb)\) 对一切 \(a \in A\) 、\(b \in B\) 、\(g \in G\) 成立。

定义 1. --- 称一个映射 \(\Phi\) 为从积 \(\mathbf{E} \times \mathbf{F}\) 到 \(\mathbf{G}\) 中的双线性映射，如果它满足下列条件：

\begin{enumerate}
\def\labelenumi{(\arabic{enumi})}
\tightlist
\item
  \(\Phi(x + x', y) = \Phi(x, y) + \Phi(x', y)\)
\end{enumerate}

对一切 \(x \in \mathbf{E}\) 、\(x' \in \mathbf{E}\) 、\(y \in \mathbf{F}\) ；

\begin{enumerate}
\def\labelenumi{(\arabic{enumi})}
\setcounter{enumi}{1}
\tightlist
\item
  \(\Phi(x, y + y') = \Phi(x, y) + \Phi(x, y')\)
\end{enumerate}

对一切 \(x \in \mathbf{E}, y \in \mathbf{F}, y' \in \mathbf{F}\) ；

\begin{enumerate}
\def\labelenumi{(\arabic{enumi})}
\setcounter{enumi}{2}
\item
  \(\Phi(ax, y) = a\Phi(x, y)\) 对一切 \(a \in \mathbf{A}\) 、\(x \in \mathbf{E}\) 、\(y \in \mathbf{F}\) ；
\item
  \(\Phi(x, yb) = \Phi(x, y)b\) 对一切 \(x \in \mathbf{E}, y \in \mathbf{F}, b \in \mathbf{B}\) 。
\end{enumerate}

张量积 \(E \otimes_{Z} F\) 上典范地带有由 \(a(x \otimes y)b = ax \otimes yb\) 刻画的一个 (A, B)-双模结构（第 III 章，第 2 版，附录 II，第 3 小节），而双线性映射 \(\Phi\) （从 \(E \times F\) 到 G 中）的数据，等价于映射 \(\Psi\) （从 \(E \otimes_{Z} F\) 到 G 中，对于 (A, B)-双模结构而言是同态，且满足 \(\Psi(x \otimes y) = \Phi(x, y)\) ，对一切 \(x \in E\) 与 \(y \in F\)）的数据。

定义 1 加于 \(\Phi\) 的条件表明，偏映射 \(d_{\Phi}(y): x \to \Phi(x, y)\) 与 \(s_{\Phi}(x): y \to \Phi(x, y)\) 分别是 E 到 G 的 A-线性映射与 F 到 G 的 B-线性映射。给交换群 \(\mathcal{L}_{\mathbf{A}}(\mathbf{E}, \mathbf{G})\) （相应地 \(\mathcal{L}_{\mathbf{B}}(\mathbf{F}, \mathbf{G})\) ）配上由 \(ub(x) = u(x). b\) （ \(u \in \mathcal{L}_{\mathbf{A}}(\mathbf{E}, \mathbf{G}), x \in \mathbf{E}, b \in \mathbf{B}\) ）（相应地 \(av(y) = a. v(y)\) （ \(a \in \mathbf{A}, v \in \mathcal{L}_{\mathbf{B}}(\mathbf{F}, \mathbf{G}), y \in \mathbf{F}\) ））所定义的右 B-模（相应地左 A-模）结构。于是，条件 (1) 至 (4) 分别等价于：

\[
s _ {\Phi} (x + x ^ {\prime}) = s _ {\Phi} (x) + s _ {\Phi} (x ^ {\prime})\tag{1'}
\]

\[
d _ {\Phi} (y + y ^ {\prime}) = d _ {\Phi} (y) + d _ {\Phi} (y ^ {\prime})\tag{2'}
\]

\[
s _ {\Phi} (a x) = a. s _ {\Phi} (x)\tag{3'}
\]

\[
d _ {\Phi} (y b) = d _ {\Phi} (y). b,\tag{4'}
\]

对 E 中的一切 \(x, x'\) 、F 中的一切 \(y, y'\) 、\(a \in A\) 、\(b \in B\) ；换言之，映射 \(d_{\Phi}\) （从 F 到 \(\mathcal{L}_{A}(E, G)\) 中）是 B-线性的，而映射 \(s_{\Phi}\) （从 E 到 \(\mathcal{L}_{B}(F, G)\) 中）是 A-线性的。按定义，我们有

\begin{enumerate}
\def\labelenumi{(\arabic{enumi})}
\setcounter{enumi}{4}
\tightlist
\item
  \(\Phi(x, y) = d_{\Phi}(y)(x) = s_{\Phi}(x)(y)\) 对一切 \(x \in \mathbf{E}, y \in \mathbf{F}\) 。
\end{enumerate}

定义 2. --- 给定 E × F 到 G 的双线性映射 \(\Phi\)，映射 \(d_{\Phi}\) （从 F 到 \(\mathfrak{L}_{\mathrm{A}}(\mathrm{E}, \mathrm{G})\) 中；相应地，映射 \(s_{\Phi}\) 从 E 到 \(\mathfrak{L}_{\mathrm{B}}(\mathrm{F}, \mathrm{G})\) 中）由 (5) 刻画，称为 \(\Phi\) 的右相伴线性映射（相应地左相伴线性映射）。

反之，B-线性映射 \(d\) （从 F 到 \(\mathfrak{L}_{\mathbf{A}}(\mathbf{E},\mathbf{G})\) 中；相应地，A-线性映射 \(s\) 从 E 到 \(\mathfrak{L}_{\mathbf{B}}(\mathbf{F},\mathbf{G})\) 中）的数据，借助公式

\[
\Phi (x, y) = d (y) (x) \quad (\text { resp. } \Phi (x, y) = s (x) (y))
\]

唯一确定一个双线性映射 \(\Phi\)，它映 \(\mathbf{E} \times \mathbf{F}\) 到 \(\mathbf{G}\) 中，且 \(d\) （相应地 \(s\)）是它的右相伴线性映射（相应地左相伴线性映射）。

定义 3. --- 称 E × F 到 G 的双线性映射 \(\Phi\) 为右退化的（相应地左退化的），如果存在 F 的一个非零元素 \(y_0\) （相应地 E 的一个非零元素 \(x_0\)），使得 \(\Phi(x, y_0) = 0\) 对一切 \(x \in \mathbf{E}\) 成立（相应地 \(\Phi(x_0, y) = 0\) 对一切 \(y \in \mathbf{F}\) 成立）。称 \(\Phi\) 为退化的，如果它是右退化的或左退化的。

\(\Phi\) 右非退化（相应地左非退化）必须且只需 \(\Phi\) 的右相伴线性映射（相应地左相伴线性映射）是单射；因此，说 \(\Phi\) 非退化，就是说相伴线性映射 \(d_{\Phi}\) 与 \(s_{\Phi}\) 都是单射。

设 \((e_i)_{i\in I}\) 与 \((f_k)_{k\in K}\) 是 E 与 F 的元素的两个族，并设 \((a_i)_{i\in I}\) 与 \((b_k)_{k\in K}\) 是 A 与 B 的元素的两个族，其中除有限个外均为零。由等式 (1) 至 (4)，对非零系数的个数作归纳，可得

\[
\Phi (\sum_ {i} a _ {i} e _ {i}, \sum_ {k} f _ {k} b _ {k}) = \sum_ {i, k} a _ {i} \Phi (e _ {i}, f _ {k}) b _ {k}.\tag{6}
\]

若 \((e_{i})\) 与 \((f_{k})\) 是模 E 与 F 的生成系，则 \(\Phi\) 由元素 \(g_{ik} = \Phi(e_{i}, f_{k})\) 完全确定。若 \((e_{i})\) 与 \((f_{k})\) 是 E 与 F 的基，且给定了 G 的元素 \(g_{ik}\) \((i \in I, k \in K)\)，则公式

\[
\Phi (\sum_ {i} a _ {i} e _ {i}, \sum_ {k} f _ {k} b _ {k}) = \sum_ {i, k} a _ {i} g _ {i k} b _ {k}\tag{6'}
\]

定义了 E × F 到 G 的一个映射，它是双线性的，且满足 \(\Phi(e_{i}, f_{k}) = g_{ik}\) 。当 \((e_{i})\) 与 \((f_{k})\) 为有限基时，称 \((\Phi(e_{i}, f_{k}))\) 为 \(\Phi\) 关于这些基的矩阵。

E × F 到 G 中的双线性映射显然构成 E × F 到 G 的映射的加法群的一个子群。另一方面，设 a（相应地 b）为 A（相应地 B）的中心的一个元素；则映射 \(a\Phi b\) （从 E × F 到 G 中，由 \((a\Phi b)(x, y) = a \cdot \Phi(x, y) \cdot b\) 定义）是双线性的。于是，E × F 到 G 中的双线性映射的集合便带有 A 与 B 的中心上的一个双模结构。

设 E′（相应地 F′）为左 A-模（相应地右 B-模），u（相应地 v）为 E 到 E′（相应地 F 到 F′）的同态，\(\Phi'\) 为 \(E' \times F'\) 到 G 中的双线性映射。\(\Phi'\) （关于 u 与 v）的原像是 E × F 到 G 的双线性映射 \(\Phi\)，它定义为

\[
\Phi (x, y) = \Phi^ {\prime} (u (x), v (y)) \quad (x \in \mathrm{E}, y \in \mathrm{F}).
\]

容易验证，

\[
d _ {\Phi} (y) = d _ {\Phi^ {\prime}} (\nu (y)) \circ u \quad \text { and } \quad s _ {\Phi} (x) = s _ {\Phi^ {\prime}} (u (x)) \circ \nu
\]

对一切 \(x \in \mathbf{E}, y \in \mathbf{F}\) 。

设 \(\Phi\) 为 \(\mathbf{E} \times \mathbf{F}\) 到 \(\mathbf{G}\) 中的双线性映射，\(h\) 为 \(\mathbf{G}\) 到另一个 (A, B)-双模 \(\mathbf{G}'\) 中的一个（对于 (A, B)-双模结构的）同态。则 \(h \circ \Phi\) 是 \(\mathbf{E} \times \mathbf{F}\) 到 \(\mathbf{G}'\) 中的双线性映射。

